% GENERATED FILE. Edit canonical lessons, then run npm run generate:books.
% canonical-source-hash: fnv1a64:d7498a6e751428a9
% canonical-lessons: GE-C497-trotzig, GE-C497-schimpfen, GE-C497-loben, GE-C497-necken, GE-C497-taschengeld, GE-R497-review-490-497

\chapter{Frech und trotzig}
\label{ch:ge-frech-und-trotzig}
\hlchaptermodality{497}

\begin{chapteropening}
\textbf{By the end of this chapter:} \emph{I can talk about children who misbehave.}

Complete the last lesson of chapter 497: I can talk about children who misbehave.
\end{chapteropening}

\section[trotzig]{trotzig — defiant}

\label{lesson:GE-C497-trotzig}



\begin{quote}
Before the new one: say the German for pram, then the German for dummy.
\end{quote}

\subsection*{You'll want to know: trotzig}
\textbf{trotzig} — \textquotedblleft{}defiant\textquotedblright{}.

\begin{quote}
Der Kleine ist müde und deshalb trotzig.
\end{quote}

\begin{quote}
\emph{The little one is tired and so he's being defiant.}
\end{quote}

\emph{Die Trotzphase} — the terrible twos.

\begin{cousinweb}
\emph{trotzig} is built on \emph{Trotz} \textquotedblleft{}defiance\textquotedblright{}, from Middle High German \emph{traz} (EtymDB), with \emph{-ig}.
\end{cousinweb}

\begin{grammarlens}[title={What you've built}]
One more word for telling, explaining and describing at B1.
\end{grammarlens}

\subsection*{Your turn}
\begin{itemize}
\raggedright
  \item \emph{Say it:} \emph{trotzig}
  \item \emph{Say it:} \emph{trotzig}, once more
  \item \emph{Say it:} a sentence of your own with \emph{trotzig}
  \item \emph{Recall:} say the German for pram, then the German for dummy, then say \emph{trotzig} again
\end{itemize}

\begin{tcolorbox}[breakable,colback=teal!4,colframe=teal!35!black,title={Before you move on}]
What does trotzig mean? (\textquotedblleft{}Defiant\textquotedblright{}.) Say it once more.
\end{tcolorbox}
\section[schimpfen]{schimpfen — to scold; to grumble}

\label{lesson:GE-C497-schimpfen}



\begin{quote}
Before the new one: say the German for dummy, then the German for defiant.
\end{quote}

\subsection*{You'll want to know: schimpfen}
\textbf{schimpfen} — \textquotedblleft{}to scold; to grumble\textquotedblright{}.

\begin{quote}
Die Mutter schimpft, weil die Kinder so laut sind.
\end{quote}

\begin{quote}
\emph{The mother is scolding because the children are so noisy.}
\end{quote}

\emph{Mit jemandem schimpfen} — to tell someone off. \emph{Über etwas schimpfen} — to grumble about something.

\begin{cousinweb}
\emph{schimpfen} goes back to Middle High German \emph{schimpfen} \textquotedblleft{}to joke, to mock\textquotedblright{}; the meaning later shifted to scolding. EtymDB lists the German word without a source.
\end{cousinweb}

\begin{grammarlens}[title={What you've built}]
One more word for telling, explaining and describing at B1.
\end{grammarlens}

\subsection*{Your turn}
\begin{itemize}
\raggedright
  \item \emph{Say it:} \emph{schimpfen}
  \item \emph{Say it:} \emph{schimpfen}, once more
  \item \emph{Say it:} a sentence of your own with \emph{schimpfen}
  \item \emph{Recall:} say the German for dummy, then the German for defiant, then say \emph{schimpfen} again
\end{itemize}

\begin{tcolorbox}[breakable,colback=teal!4,colframe=teal!35!black,title={Before you move on}]
What does schimpfen mean? (\textquotedblleft{}To scold; to grumble\textquotedblright{}.) Say it once more.
\end{tcolorbox}
\section[loben]{loben — to praise}

\label{lesson:GE-C497-loben}



\begin{quote}
Before the new one: say the German for defiant, then the German for to scold.
\end{quote}

\subsection*{You'll want to know: loben}
\textbf{loben} — \textquotedblleft{}to praise\textquotedblright{}.

\begin{quote}
Die Lehrerin lobt den Jungen für seine Arbeit.
\end{quote}

\begin{quote}
\emph{The teacher praises the boy for his work.}
\end{quote}

The noun: \emph{das Lob} — praise.

\begin{cousinweb}
\emph{loben} is inherited from Middle High German \emph{loben}, from Proto-Germanic \emph{lubōną} \textquotedblleft{}to praise\textquotedblright{} (EtymDB).
\end{cousinweb}

\begin{grammarlens}[title={What you've built}]
One more word for telling, explaining and describing at B1.
\end{grammarlens}

\subsection*{Your turn}
\begin{itemize}
\raggedright
  \item \emph{Say it:} \emph{loben}
  \item \emph{Say it:} \emph{loben}, once more
  \item \emph{Say it:} a sentence of your own with \emph{loben}
  \item \emph{Recall:} say the German for defiant, then the German for to scold, then say \emph{loben} again
\end{itemize}

\begin{tcolorbox}[breakable,colback=teal!4,colframe=teal!35!black,title={Before you move on}]
What does loben mean? (\textquotedblleft{}To praise\textquotedblright{}.) Say it once more.
\end{tcolorbox}
\section[necken]{necken — to tease}

\label{lesson:GE-C497-necken}



\begin{quote}
Before the new one: say the German for to scold, then the German for to praise.
\end{quote}

\subsection*{You'll want to know: necken}
\textbf{necken} — \textquotedblleft{}to tease\textquotedblright{}.

\begin{quote}
Mein großer Bruder neckt mich gern wegen meiner Haare.
\end{quote}

\begin{quote}
\emph{My big brother likes to tease me about my hair.}
\end{quote}

Gentle, friendly teasing.

\begin{cousinweb}
The origin of \emph{necken} is not settled; EtymDB lists the German word without a source.
\end{cousinweb}

\begin{grammarlens}[title={What you've built}]
One more word for telling, explaining and describing at B1.
\end{grammarlens}

\subsection*{Your turn}
\begin{itemize}
\raggedright
  \item \emph{Say it:} \emph{necken}
  \item \emph{Say it:} \emph{necken}, once more
  \item \emph{Say it:} a sentence of your own with \emph{necken}
  \item \emph{Recall:} say the German for to scold, then the German for to praise, then say \emph{necken} again
\end{itemize}

\begin{tcolorbox}[breakable,colback=teal!4,colframe=teal!35!black,title={Before you move on}]
What does necken mean? (\textquotedblleft{}To tease\textquotedblright{}.) Say it once more.
\end{tcolorbox}
\section[das Taschengeld]{das Taschengeld — pocket money}

\label{lesson:GE-C497-taschengeld}



\begin{quote}
Before the new one: say the German for to praise, then the German for to tease.
\end{quote}

\subsection*{You'll want to know: das Taschengeld}
\textbf{das Taschengeld} — \textquotedblleft{}pocket money\textquotedblright{}.

\begin{quote}
Von seinem Taschengeld kauft er Comics.
\end{quote}

\begin{quote}
\emph{He buys comics with his pocket money.}
\end{quote}

No plural.

\begin{cousinweb}
\emph{Taschengeld} is built on \emph{Tasche} \textquotedblleft{}pocket\textquotedblright{}, inherited from Old High German \emph{tasca}, and \emph{Geld} \textquotedblleft{}money\textquotedblright{}, from Old High German \emph{gelt} \textquotedblleft{}payment\textquotedblright{} (EtymDB).
\end{cousinweb}

\begin{grammarlens}[title={What you've built}]
One more word for telling, explaining and describing at B1.
\end{grammarlens}

\subsection*{Your turn}
\begin{itemize}
\raggedright
  \item \emph{Say it:} \emph{das Taschengeld}
  \item \emph{Say it:} \emph{das Taschengeld}, once more
  \item \emph{Say it:} a sentence of your own with \emph{das Taschengeld}
  \item \emph{Recall:} say the German for to praise, then the German for to tease, then say \emph{das Taschengeld} again
\end{itemize}

\begin{tcolorbox}[breakable,colback=teal!4,colframe=teal!35!black,title={Before you move on}]
What does das Taschengeld mean? (\textquotedblleft{}Pocket money\textquotedblright{}.) Say it once more.
\end{tcolorbox}
\section[(dialogue)]{Words from chapters 490–497 — a review}

\label{lesson:GE-R497-review-490-497}



\begin{quote}
Say the German for father-in-law and stepfather.
\end{quote}

\subsection*{Your turn}
Say these aloud:

\begin{itemize}
\raggedright
  \item father-in-law, stepfather, great-grandfather, twin, only child
  \item half-brother, godfather, godchild, spouse, offspring
  \item to fall in love, to flirt, to cuddle, to kiss, kiss
  \item affection, hug, tender, feeling of safety, solidarity
  \item to squabble, to sulk, to get along, to reconcile, reconciliation
  \item to bring up, upbringing, to spoil, to obey, puberty
  \item to breastfeed, nappy, to change, pram, dummy
  \item defiant, to scold, to praise, to tease, pocket money
\end{itemize}

\begin{tcolorbox}[breakable,colback=teal!4,colframe=teal!35!black,title={Before you move on}]
Father-in-law? (\textbf{der Schwiegervater}.) Stepfather? (\textbf{der Stiefvater}.) Great-grandfather? (\textbf{der Urgroßvater}.) Twin? (\textbf{der Zwilling}.) Only child? (\textbf{das Einzelkind}.) Half-brother? (\textbf{der Halbbruder}.) Godfather? (\textbf{der Pate}.) Godchild? (\textbf{das Patenkind}.) Spouse? (\textbf{der Ehepartner}.) Offspring? (\textbf{der Nachwuchs}.) To fall in love? (\textbf{verlieben}.) To flirt? (\textbf{flirten}.) To cuddle? (\textbf{kuscheln}.) To kiss? (\textbf{küssen}.) Kiss? (\textbf{der Kuss}.) Affection? (\textbf{die Zuneigung}.) Hug? (\textbf{die Umarmung}.) Tender? (\textbf{zärtlich}.) Feeling of safety? (\textbf{die Geborgenheit}.) Solidarity? (\textbf{der Zusammenhalt}.) To squabble? (\textbf{zanken}.) To sulk? (\textbf{schmollen}.) To get along? (\textbf{vertragen}.) To reconcile? (\textbf{versöhnen}.) Reconciliation? (\textbf{die Versöhnung}.) To bring up? (\textbf{erziehen}.) Upbringing? (\textbf{die Erziehung}.) To spoil? (\textbf{verwöhnen}.) To obey? (\textbf{gehorchen}.) Puberty? (\textbf{die Pubertät}.) To breastfeed? (\textbf{stillen}.) Nappy? (\textbf{die Windel}.) To change? (\textbf{wickeln}.) Pram? (\textbf{der Kinderwagen}.) Dummy? (\textbf{der Schnuller}.) Defiant? (\textbf{trotzig}.) To scold? (\textbf{schimpfen}.) To praise? (\textbf{loben}.) To tease? (\textbf{necken}.) Pocket money? (\textbf{das Taschengeld}.)
\end{tcolorbox}
